% Run GCT_KRON_20261120; GPT-6 Astra Pro; effort not exposed.
\section{An actual multiplicity obstruction for the unpadded permanent}

\subsection*{A-030. The Hessian bracket and its unique ambient copy}
Work over \(\mathbb C\). Let \(n\geq3\), \(R=n^2\),
\(W=\operatorname{Mat}_n(\mathbb C)\), and put
\[
 q=R(n-2),\qquad d=R+1,\qquad
 \lambda=(n+q+1,3,\underbrace{2,\ldots,2}_{R-2}).       \tag{A-030.1}
\]
Thus \(|\lambda|=nd\) and \(\ell(\lambda)=R\).
In the notation of the main report, with the permanent unpadded
and of the same size as the determinant, one has
\[
 a_\lambda(d[n])=1,\qquad
 K_n(\lambda)=0,\qquad
 q_\lambda^{\,n}(d)=1.                               \tag{A-030.2}
\]
In particular this is an actual multiplicity obstruction
\(q_\lambda^{\,n}(d)>K_n(\lambda)\), for every \(n\geq3\).
For example, \(n=3\) gives \(d=10\) and
\(\lambda=(13,3,2,2,2,2,2,2,2)\).
It is outside the large-padding range of the B\"urgisser--Ikenmeyer--Panova
theorem. No assertion of historical priority is made.
In fact this is an ordinary rectangular Kronecker occurrence
obstruction as well:
\[
               g(\lambda,(d^n),(d^n))=0.                       \tag{A-030.2a}
\]
Accordingly the symmetric Kronecker multiplicity also vanishes;
there is no orbit-to-closure multiplicity loss in this family.

\paragraph{Known inputs and the additional calculation.}
The permanent Hessian used below is exactly the classical
Mignon--Ressayre point, not a newly discovered full-rank point:
\emph{A quadratic bound for the determinant and permanent problem},
author preprint dated July 20, 2004, Section 3.2.2,
Proposition 3.1, printed p.~8. Their Lemma 3.4 on the same page
gives its first and second derivatives, and pp.~9--10 prove
invertibility. Their Theorem 2, p.~10, is the resulting quadratic
lower bound on determinantal complexity. We use Proposition 3.1
as a known result. The determinant's Hessian power identity is
classical as well: Thomas Muir, \emph{The Hessian of a General
Determinant}, \emph{Transactions of the Royal Society of Edinburgh}
40 (1905), 203--207, DOI \texttt{10.1017/S008045680003427X};
the publisher's extract explicitly states the power
\(n(n-2)\). The short computation below fixes the nonzero
constant for ordinary second derivatives and a specified basis.

The new work in this report is the explicit assembly of the Hessian
bracket, its Schur type, and F-024's ambient multiplicity-one
calculation into (A-030.2), together with an exact determinant
normalization check. This description does not claim that this
assembly is absent from the literature. Hessian equations for
determinant orbit closures and a stronger quadratic border bound
are already established by Landsberg--Manivel--Ressayre,
\emph{Hypersurfaces with degenerate duals and the Geometric
Complexity Theory Program},
\emph{Commentarii Mathematici Helvetici} 88 (2013), 469--484,
Theorem 1.1.1, printed p.~470; Theorem 2.3.1, p.~474;
and Theorem 3.1.1, p.~476. Their Section 3.2, p.~476,
explicitly distinguishes one ideal copy from an entire isotypic
component. Here F-024 supplies precisely the ambient uniqueness
that turns the explicit ideal copy into the full isotypic component.

\paragraph{The Hessian power identity, with normalization.}
Use the ordered matrix-entry basis on \(W\). For a homogeneous
degree-\(n\) polynomial \(F\) set
\[
 \mathcal H(F)(x)=\det\!\left(
           \frac{\partial^2 F}{\partial x_i\partial x_j}(x)
                         \right)_{1\leq i,j\leq R}.
\]
This is a homogeneous form of degree \(q\) in \(x\), and a
polynomial of degree \(R\) in the coefficients of \(F\).
For \(D=\det_n\), the classical identity, in this normalization, is
\[
 \mathcal H(D)=c_nD^{n(n-2)},\qquad
 c_n=(n-1)(-1)^{n(n+1)/2-1}\ne0.                     \tag{A-030.3}
\]
Indeed, the second differential at the identity is
\(\operatorname{tr}(X)\operatorname{tr}(Y)-\operatorname{tr}(XY)\).
Relative to the entrywise pairing its operator is
\(X\mapsto\operatorname{tr}(X)I-X^{\mathsf T}\).
Its eigenvalues are \(n-1\) on the scalar line, \(-1\) on the
symmetric trace-zero matrices, and \(1\) on the skew-symmetric
matrices. This gives \(c_n\). For an invertible matrix \(A\),
the change \(X\mapsto A^{-1}X\) has determinant
\(\det(A)^{-n}\), while the second differential gains the scalar
\(\det(A)\). Taking its \(R\)-by-\(R\) determinant gives
\(c_n\det(A)^{R-2n}\). Polynomial continuation proves (A-030.3).
This verifies the constant rather than imposing a convention from
a differently normalized polarization.

\paragraph{A polynomial covariant in one irreducible module.}
The action is \((g\cdot F)(x)=F(g^{-1}x)\), so the chain rule gives
\[
 \mathcal H(g\cdot F)(x)
       =\det(g)^{-2}\mathcal H(F)(g^{-1}x).             \tag{A-030.4}
\]
Consequently \(\mathcal H\) is intrinsically valued in
\((\det W^*)^{\otimes2}\otimes\operatorname{Sym}^{q}W^*\).
For \(a=n+q-1\), form the one-form covariant
\[
 \Theta(F)=q\mathcal H(F)\,\mathrm dF
                         -nF\,\mathrm d\mathcal H(F).          \tag{A-030.5}
\]
Its degree in the coefficients of \(F\) is \(R+1=d\).
Its polynomial coefficients have degree \(a\) in \(x\).
Euler's identity gives
\[
                       \Theta(F)_x(x)=0.                      \tag{A-030.6}
\]
The kernel of multiplication
\(\operatorname{Sym}^aW^*\otimes W^*
\longrightarrow\operatorname{Sym}^{a+1}W^*\)
is \(\mathbb S_{(a,1)}W^*\), by the characteristic-zero
Pieri decomposition. Thus (A-030.5) is a polynomial
\(\operatorname{GL}(W)\)-equivariant map to the irreducible module
\[
 U=(\det W^*)^{\otimes2}\otimes\mathbb S_{(a,1)}W^*
                 =\mathbb S_\lambda W^*.                       \tag{A-030.7}
\]
The two determinant factors add two boxes to every one of the
\(R\) rows. They do not produce a two-row representation:
the resulting highest-weight partition is exactly (A-030.1).

Equation (A-030.3), together with \(q=n\,n(n-2)\), implies
\(\Theta(D)=0\). Equivariance then gives zero on the determinant
orbit. Since every coefficient of \(\Theta\) is a polynomial
in the coefficients of \(F\), it gives zero on the full affine
orbit closure \(\Omega_n\) as well. This is a closure assertion,
not only an orbit-stabilizer calculation.

\paragraph{Nonvanishing on the permanent.}
Let \(P=\operatorname{perm}_n\) and let \(A\) have all entries
equal to one except \(A_{11}=1-n\). Mignon--Ressayre,
Proposition 3.1 and Lemma 3.4 at the locators above, give
\[
 P(A)=0,\qquad
 \frac{\partial P}{\partial x_{11}}(A)=(n-1)!,\qquad
 \operatorname{rank}\operatorname{Hess}(P)(A)=n^2.               \tag{A-030.8}
\]
Thus \(\mathcal H(P)(A)\ne0\), and evaluating (A-030.5) yields
\[
 \Theta(P)_A=q\mathcal H(P)(A)\,\mathrm dP_A\ne0.                \tag{A-030.9}
\]
In particular the coefficient map
\[
 \Theta^*:U^*\longrightarrow
       \mathbb C[\operatorname{Sym}^nW^*]_d
       =\operatorname{Sym}^{d}(\operatorname{Sym}^nW)            \tag{A-030.10}
\]
is nonzero. Its source \(U^*=\mathbb S_\lambda W\) is
irreducible, so it is injective. Its image is contained in the
ideal of \(\Omega_n\), whereas its restriction to the permanent
orbit closure is nonzero, hence injective on this simple module.

Now F-024 gives, for every outer degree \(d\geq3\) and inner
degree \(n\geq3\),
\[
 a_{(d(n-2)+3,\,3,\,2^{d-3})}(d[n])=1
 \quad\hbox{when the ambient dimension is at least }d-1.         \tag{A-030.11}
\]
Here \(d=R+1\) and \(\dim W=R=d-1\), and the first part
\(d(n-2)+3\) is \(n+q+1\). Therefore the copy in
(A-030.10) is the entire ambient isotypic component. It dies
on the determinant closure and survives on the permanent closure.
Complete reducibility in each finite degree proves (A-030.2).
In particular it is the ambient uniqueness, in addition to the
separating equation, that proves the exact multiplicity values by
this argument. The following independent vanishing calculation
gives a second route to the determinant side.

\paragraph{The ordinary Kronecker coefficient already vanishes.}
Write \(W=A_0\otimes B_0\) with
\(\dim A_0=\dim B_0=n\), and put
\[
 E=n+q=n(n-1)^2,\qquad \nu=(E-1,1),\qquad
                         t=d-2n=(n-1)^2.
\]
As a \(\operatorname{GL}(W)\)-module,
\(\mathbb S_\lambda W=(\det W)^2\otimes\mathbb S_\nu W\).
On restriction to \(\operatorname{GL}(A_0)\times
\operatorname{GL}(B_0)\),
\(\det(A_0\otimes B_0)^2=(\det A_0)^{2n}(\det B_0)^{2n}\).
The Schur--Weyl definition of Kronecker multiplicities therefore gives
\[
 g(\lambda,(d^n),(d^n))
                         =g((E-1,1),(t^n),(t^n)).              \tag{A-030.14}
\]
The standard hook \((E-1,1)\) has multiplicity equal to the
number of removable corners of \((t^n)\) minus one in the
tensor square of its Specht module. This is zero for a rectangle;
equivalently it is the degree-one coefficient of F-003's
rectangular-hook product. This proves (A-030.2a).
Since \(K_n(\lambda)\) is bounded by the symmetric Kronecker
multiplicity, which in turn is bounded by the ordinary one,
it independently gives \(K_n(\lambda)=0\). Thus the example
does not rely on substituting an orbit coefficient for a larger
or smaller closure coefficient.

\paragraph{A saturation hole with actual permanent positivity.}
For every integer \(s\geq1\), the same determinant twist gives
\begin{align}
 g(s\lambda,((sd)^n),((sd)^n))
   &=g((s(E-1),s),((st)^n),((st)^n))\notag\\
   &=[z^s]\prod_{j=2}^{n}(1-z^j)^{-1}.                \tag{A-030.15}
\end{align}
The second equality is the known rectangular two-row formula
F-002: it is \(p_s(n,st)-p_{s-1}(n,st)\), where \(p_u(n,st)\)
counts partitions of \(u\) in an \(n\)-by-\(st\) rectangle.
Since \(st=s(n-1)^2\geq s\), the width bound has no effect at
these two coefficients, giving the displayed product.
Consequently the unstretched coefficient is zero, the
twice-stretched coefficient is one, and the coefficient is positive
for every \(s\geq2\): all such integers are sums of twos and threes,
and \(n\geq3\).

Moreover,
\[
                         q_{s\lambda}^{\,n}(sd)>0
                         \quad\text{for every }s\geq1.          \tag{A-030.16}
\]
Indeed choose a nonzero highest-weight vector in the copy
(A-030.10) after restriction to the permanent orbit closure.
That closure is irreducible, so its coordinate algebra is a
domain. Its \(s\)-th power is a nonzero highest-weight vector
of weight \(s\lambda\) and degree \(sd\). This proves
(A-030.16). Thus the ordinary Kronecker saturation hole here
has genuine permanent positivity, rather than just a positive
ambient plethysm coefficient. Equations (A-030.15)--(A-030.16)
do not determine the determinant-closure multiplicities
\(K_n(s\lambda)\) for \(s\geq2\); positive ordinary Kronecker
multiplicity is not a boundary-extension theorem.

\paragraph{Optional exact normalization of the permanent Hessian.}
The following explicit scalar sharpens the nonsingularity input;
it is not needed in order to invoke the known Proposition 3.1:
\[
 \mathcal H(P)(A)=
 ((n-3)!)^{n^2}(-2)^{(n-2)^2}
             (n-2)^{4(n-1)}(n-1)^{2n}.                        \tag{A-030.12}
\]
Here is a check that also specifies all normalization choices.
More generally take \(X(t)\) with entry \(11\) equal to \(t\)
and every other entry equal to one. Put
\(k=n-1\), \(b=(n-2)!\),
\(c=(n-3)!(t+n-3)\). Two Hessian entries in the same row or
column of the underlying matrix have value zero. Other entries
have value \(b\) if one of their rows or columns is numbered
one, and value \(c\) otherwise. This follows either from
the permanent minor formula, or directly from Lemma 3.2 of
Mignon--Ressayre, p.~7.

Decompose \(\mathbb C^k=\mathbb C\mathbf1\oplus E\), where
\(E\) is the sum-zero subspace, using the permutations of the
last \(k\) rows and columns. On the bottom-block summand
\(E\otimes E\), of dimension \((k-1)^2\), the Hessian has
eigenvalue \(c\). The first-row and matching column-constant
sum-zero modes, and the analogous first-column modes, have
two-by-two blocks
\[
 \begin{pmatrix}0&-b\sqrt{k}\\-b\sqrt{k}&-c(k-1)\end{pmatrix},
\]
of determinant \(-b^2k\), repeated \(2(k-1)\) times in total.
On the remaining four-dimensional constant subspace, use the
unit vectors of the \(11\) coordinate, first row, first column,
and bottom block. The matrix is
\[
 \begin{pmatrix}
 0&0&0&bk\\
 0&0&bk&b\sqrt{k}(k-1)\\
 0&bk&0&b\sqrt{k}(k-1)\\
 bk&b\sqrt{k}(k-1)&b\sqrt{k}(k-1)&c(k-1)^2
 \end{pmatrix},
\]
whose determinant is \(b^4k^4\). These are an orthonormal
decomposition of the entrywise bilinear form, so their product is
\[
          \mathcal H(P)(X(t))
               =b^{4k}k^{2k+2}c^{(k-1)^2}.                     \tag{A-030.13}
\]
At \(t=-k=1-n\), \(c=-2(n-3)!\), and (A-030.12) follows.
The factorial exponent is
\(4(n-1)+(n-2)^2=n^2\). No choice of square root affects the
block determinants. Exact permutation-sum differentiation and
integer determinants for \(n=3,4,5,6\) are recorded separately;
these finite checks supplement the displayed general derivation.

\paragraph{Padding limitation.}
For the same determinant size \(n\), the report's internally
padded permanent \(x_{11}^{n-m}\operatorname{perm}_m\),
with \(m<n\), uses at most \(m^2<n^2\) essential variables.
The present partition has
\(n^2\) nonzero parts. Therefore its entire ambient isotypic
component vanishes on this subspace variety: equivalently,
\(q_\lambda^{\,m}(d)=0\), by the standard Schur length bound
for the polynomial coordinate ring on forms in fewer variables.
One can also see the failure of this particular equation directly:
the full \(R\)-variable Hessian determinant, and hence
\(\Theta\), is identically zero on such a form. Thus (A-030.2)
does not transfer to \(n>m\), and supplies no contradiction to
the large-padding no-occurrence theorem or a superpolynomial
determinantal-complexity lower bound.
For the necessary second-row rule established in F-022 and
M-016, namely \(n\geq m^c\Rightarrow\lambda_2\geq c+2\)
for an actual obstruction, the present family attains
\(\lambda_2=3\) at \(c=1\) and \(n=m\). It proves sharpness
of that endpoint only, with no conclusion about the endpoints
at \(c\geq2\).

\paragraph{Grothendieck dictionary and exact scope.}
The construction is a polynomial map between finite-dimensional
representations, and its coefficients define a homogeneous
\(\operatorname{GL}(W)\)-stable ideal in the coordinate algebra
of the affine space \(\operatorname{Sym}^nW^*\). The determinant
closure is the scheme-theoretic image closure; vanishing of the
coefficient map on its dense orbit gives a factorization through
its ideal. Schur operations here are actual characteristic-zero
representations, the setting of the SGA 6, VI, Theorem 3.3
dictionary already used in the report. After taking the degree-\(d\)
piece, (A-030.10) is an explicit kernel inclusion for the
restriction map to the determinant closure, and (A-030.9)
proves its nonzero image under restriction to the permanent
closure. F-024 computes the ambient rank of this isotypic part.
Thus this example performs the ideal-kernel and quotient-rank
comparison that positivity of a virtual class alone cannot give.
No claim of invariants commuting with arbitrary integral base
change, or of a statement in positive characteristic, is made.
